\begin{titlepage}
\centering
\vspace*{2.2cm}

{\LARGE\bfseries Predicting Structural Column Arrangements\\[0.35em]
from Architectural Floor Plans\par}

\vspace{1.1cm}
{\large A Controlled Comparison of an Algorithmic Score, a Feature-Based
Classifier, an Image-Based Classifier and a Quantum Classifier\par}

\vspace{2.6cm}
{\large MSc Thesis\par}
\vspace{0.8cm}
{\large Kritvi Gera\par}
\vspace{0.2cm}
{\large CID: 06053693\par}
\vspace{0.2cm}
{\large MSc. General Structural Engineering with Data Science\par}
\vspace{0.4cm}
{\large Imperial College London\par}

\vspace{2.8cm}
\begin{minipage}{0.8\textwidth}
\centering\small
\textit{Given an architectural floor plan, this work enumerates candidate
structural column arrangements, ranks them, and measures how closely the
highest-ranked arrangement matches the one the real building uses. Four
ranking methods are compared on one task, one corpus and one set of
labels, with the algorithm-only method retained as the baseline.}
\end{minipage}

\vfill
{\normalsize August 2026\par}
\end{titlepage}

\pagenumbering{roman}

\chapter*{Abstract}
\addcontentsline{toc}{chapter}{Abstract}
\markboth{Abstract}{}

At the start of a building's design someone must decide where the columns
will stand. The architectural drawing records the walls and the rooms; it
does not record the structure. This thesis asks whether a first proposal
for that decision can be generated and, more importantly, whether it can
be \emph{measured}.

Three public floor-plan datasets are converted into a single corpus of
29,279 plans. Because no drawing records a column layout, a label is
derived by a stated procedure: walls thinner than 160\,mm are discarded as
partitions, a grid is built from the survivors, and a column is placed at
every grid crossing that lies physically inside one of those walls. The
resulting arrangement is an equivalent-frame idealisation of a
wall-bearing structure rather than a recovered design, and the thesis is
explicit that every subsequent figure is conditional on that procedure.
The single arbitrary constant in it, the 160\,mm threshold, is varied
between 140 and 200\,mm and the effect on the headline result is reported.

Candidate arrangements are enumerated from the parsed grid and ranked by a
dimensionless geometric score built from five measures: column density, a
load-free index of bending demand, grid irregularity, the fraction of
columns standing away from walls, and the repetition of bay sizes. The
score contains no member sizes, no material prices and no load values. Its
five weights are obtained two independent ways, once by equal weighting
and once by learning them from 52,724 paired comparisons in which the real
arrangement of a plan must outrank alternatives generated on that same
plan's grid.

Four ranking variants are then compared. The baseline uses the score
alone. The other three add a learned classifier that re-orders the
score's fifteen best candidates: a gradient-boosted model reading six
scale-independent measurements, a convolutional network reading 64 by 64
pictures of the arrangement, and a four-qubit variational quantum
classifier. All three are trained on identical examples and labels, with
implausible examples produced by moving column positions rather than by
editing the numbers a model reads, which removes a circularity present in
an earlier version of this work and reduces the feature classifier's
reported score from 0.947 to 0.862.

Three findings are reported. First, the score weights matter far more than
the choice of classifier: learned weights move the median rank of the
correct arrangement, among roughly four thousand candidates, from about
2,400 to nine, whereas changing classifier moves it by one or two places.
Second, the image classifier's advantage over the feature classifier
(0.995 against 0.860 ROC AUC) is attributable to information rather than
architecture; blanking the wall layer from its input returns it to 0.874,
level with the feature model. Third, the four-qubit circuit matches its
capacity-matched classical control at every training-set size and gains
only 0.007 ROC AUC from fifty times more data, while equally informed tree
ensembles gain 0.075, which locates its ceiling in expressive capacity
rather than in data volume.

No variant separates from the baseline at conventional significance on the
seven held-out plans. That negative result is reported as the outcome of
the experiment, together with the measures taken to make it trustworthy:
candidate coverage is stated before any accuracy figure, splits are
grouped by plan and repeated, and every hand-chosen constant is varied and
its influence measured.

\cleardoublepage
\tableofcontents
\cleardoublepage
\listoftables
\addcontentsline{toc}{chapter}{List of Tables}
\cleardoublepage
\listoffigures
\addcontentsline{toc}{chapter}{List of Figures}

\chapter*{Nomenclature}
\addcontentsline{toc}{chapter}{Nomenclature}
\markboth{Nomenclature}{}

The report avoids specialist vocabulary where plain wording will serve.
The terms below are unavoidable and are used consistently throughout.

\begin{description}[style=nextline,leftmargin=3.2cm,labelwidth=3cm]
\item[Column arrangement] The set of positions at which columns stand on
one floor of a building.
\item[Load-bearing wall] A wall that carries the weight of the structure
above it, as opposed to a partition that merely divides rooms.
\item[Grid line] A straight reference line running the full width or depth
of a plan. Columns are placed where two grid lines cross.
\item[Span] The distance between two adjacent columns; equivalently the
width of one structural bay.
\item[Label, or ground truth] The arrangement treated as correct when
accuracy is measured. Architectural drawings record none, so it is derived
from the load-bearing walls by the procedure of Chapter~3.
\item[Judge] A machine-learning model that estimates how plausible a
candidate arrangement is. Three are compared, and the baseline uses none.
\item[Candidate set] Every arrangement the generator produces for one
plan, from which each variant must choose.
\item[ROC AUC] A measure of how well a model separates two classes,
running from 0.5 for a model no better than chance to 1.0 for a perfect
one.
\item[PR AUC] A companion measure that is more informative than ROC AUC
when one class is rare; both are reported throughout.
\item[Kendall's $\tau$] A rank correlation running from $-1$ to $+1$, used
here to express how far a ranking moves when a weight is perturbed.
\item[Chamfer distance] The average distance from each column in one
arrangement to the nearest column in another; a measure of near-agreement
that does not require an exact match.
\end{description}

\cleardoublepage
\pagenumbering{arabic}
